%% Section 2: Institutional Context

\section{The Naira Crunch: Institutional Context}
\label{sec:context}

\subsection{The Policy Shock}

On October~26, 2022, the Central Bank of Nigeria (\CBN) announced the
redesign of its three highest-denomination banknotes (₦200, ₦500, and
₦1,000) and set a January~31, 2023 deadline for old notes to be
deposited and replaced \citep{cbn2022circular}. Daily ATM withdrawals
were capped at ₦20,000 (approximately \$43 at the prevailing
rate) and over-the-counter withdrawals at ₦100,000 per week.
The stated rationale was to combat counterfeiting, reduce the volume
of currency held outside the banking system, and accelerate the
adoption of digital payments in support of the \CBN's broader
cashless-economy agenda.

The roll-out of new notes proved severely inadequate to demand. By
mid-January 2023, queues outside commercial banks and ATMs stretched
for hours across major Nigerian cities. Reports from the field and
media documentation consistently describe near-total unavailability
of new notes in rural areas, with POS machines dysfunctional and
markets at a standstill.\footnote{%
  The \emph{NBS Enterprise Telephone Survey}, \emph{SBM Intelligence}
  rapid surveys, and contemporaneous New Agency of Nigeria \emph{NAN} wire reports provide
  granular documentation. These sources are not used for
  identification; they motivate the timing of the shock indicators.}
The naira crunch was by this point a genuine constraint on economic
activity, not merely an inconvenience.

In February 2023 Supreme Court of Nigeria ruled that old notes must
remain legal tender until December 2023, providing partial relief.
Concurrently, the \CBN{} accelerated distribution of new currency.
By mid-2023 most commercial channels had normalised, though
\NLPS{} survey evidence (Round~11, April 2024) and the 2025 \WBES{}
both indicate persistent behavioural shifts in payment habits.

Table~\ref{tab:shock_timeline} summarises the shock timeline.

\begin{table}[H]
\begin{threeparttable}
\caption{Naira Crunch: Policy Timeline}
\label{tab:shock_timeline}
\small
\begin{tabularx}{\textwidth}{lX}
\toprule
\textbf{Date} & \textbf{Event} \\
\midrule
October 26, 2022
  & \CBN{} announces naira redesign; caps daily ATM withdrawals
    at ₦20,000; sets January 31 deadline for old notes. \\[4pt]
November--December 2022
  & Transition period: old notes recalled, new notes printed but
    distribution severely below demand. Early shortages in
    smaller cities and rural areas. \\[4pt]
January--February 2023
  & \emph{Peak crunch}: near-complete ATM dry-outs in many states;
    hour-long bank queues; POS network failures; informal market
    disruption. \\[4pt]
February 8, 2023
  & Supreme Court orders old notes remain legal tender;
    \CBN{} given three weeks to increase note supply. \\[4pt]
March--June 2023
  & Gradual recovery as new note circulation increases;
    cash availability improves progressively by state. \\[4pt]
July 2023 onward
  & Post-crunch: cash supply normalised in most states;
    some permanent shift toward digital payment habits
    documented in subsequent surveys. \\
\bottomrule
\end{tabularx}
\begin{tablenotes}
\footnotesize
\item \textit{Notes:} Timeline constructed from \CBN{} official
  circulars, Supreme Court records, and National Bureau of Statistics
  rapid-survey documentation. Shock indicator periods follow this
  timeline exactly; see Section~\ref{sec:design} for period definitions.
\end{tablenotes}
\end{threeparttable}
\end{table}

\subsection{Nigeria's Payment Landscape}
\label{sub:payment_landscape}

Understanding why fintech adoption creates differential exposure
requires appreciating Nigeria's payment infrastructure at the time of
the shock. According to the 2021 \emph{Global Findex}, approximately
45 percent of Nigerian adults held a formal bank account, well below
South Africa (85\%).
Bank branch density is heavily concentrated in Lagos, Abuja, Port
Harcourt, and Kano; many northern and rural states have fewer than
two commercial bank branches per 100,000 inhabitants \citep{IMF2022FAS}.

Against this low-baseline context, fintech adoption expanded rapidly
between 2018 and 2022. Three developments are particularly relevant.
First, mobile money operators (notably OPay and PalmPay) scaled
their agent networks aggressively, reaching peri-urban and rural
areas bypassed by traditional banking. Second, the \CBN{}'s
mandatory tier-1 Know-Your-Customer (\textsc{kyc}) relaxation in 2018
lowered barriers to opening mobile wallets. Third, \textsc{ussd} banking
(dialling a short code to transact without internet) penetrated widely
through the mobile network operator infrastructure.

A concern is whether these \CBN{}-directed policies render fintech
adoption endogenous to the 2022 crunch that if states that responded most to
the 2018 \textsc{kyc} relaxation were also anticipating a cashless policy,
the pre-shock fintech index might proxy for policy expectations rather than
exogenous infrastructure variation. Three features of the episode address
this concern. First, the specific form of the 2022 shock, a forced currency
redesign combined with withdrawal caps, was not a predictable consequence of
the 2018 \textsc{kyc} relaxation; the two episodes are legally and operationally
distinct, separated by four years. Second, the \NLPS{} Round~6 financial
service module, conducted 11~days before the announcement, directly records
\emph{realised} adoption, not plans or expectations. Also, the comparison of
Round~6 to Round~7 (February 2023) shows significant changes in financial
service usage, confirming agents had not positioned in advance. Third,
cross-state variation in fintech take-up shows differential adoption of
mobile infrastructure, driven by geography, agent-network rollout costs, and
pre-existing mobile penetration, rather than differential anticipation of
currency policy, as the geographic analysis in Section~\ref{sub:iv} shows.

By October 2022, the \NLPS{} Round~6 financial services module
(conducted 11 days before the \CBN{} announcement) recorded the
following state-level averages. 14.7 percent of households held a
commercial bank account; 86.2 percent had used mobile phone or
\textsc{ussd} banking; 5.9 percent held a mobile money account;
and 5.1 percent used a mobile banking application. The composite
fintech index, the simple average of the four digital-payment
categories, ranged from 19.8\,\% (Oyo) to 40.5\,\% (Taraba) across
states, providing the cross-sectional variation the identification exploits.

States with low fintech penetration are not simply poorer.
Appendix~Table~\ref{atab:fintech_correlates} shows that the
state-level fintech index is uncorrelated with pre-2022
nighttime light intensity ($r = -0.19$) and with pre-trend \NTL{}
growth ($r = 0.10$), supporting the parallel-trends assumption
formalised in Section~\ref{sec:design}.
